Security operations centres are asked to watch LLM agents that read untrusted content and act
through tools. We measure what a SIEM fed by the application's own telemetry can see of attacks on
such agents, with an open pack of 25 Microsoft Sentinel rules (23 also for Wazuh), 112{,}894 items from 14
public datasets, two open injection classifiers, and live attacks on twelve models from 8B
open-weight models to Claude Opus~5.5. Content signatures are precise and blind: they fire on under
1\% of benign prompts, recall 12\% of attack prompts, and lose almost everything to rewording. In this
application, against models that comply, rules over what the agent did or produced caught all 269 runs in which the attack
succeeded. Against
frontier models those rules saw 13\% of 300 attack sessions, and none on Opus~5.5, because the attacks
failed. In a pre-registered confirmation, logging provider refusals as events and matching refusal language in
the answer raised coverage to 70\% (72\% on Opus) at an 8\% benign rate: for a robust model, the evidence of an attack
is the refusal. We also show that the model, the application log and a commercial guardrail inspect
different parts of the same input, so an instruction can be followed while absent from the log; and
that a failed injection is usually relayed to the human: nine of twelve models, from small open models
to gpt-5-mini, presented an injected request to email an internal document to an outside address as
the user's next step in 61\% to 100\% of 1{,}440 runs and never warned, while the agent itself almost
never acted; only Claude Sonnet~5.5 and Opus warned. No action-based rule sees this path. We release the rules, the agent, the
pre-registrations, the raw runs and the analysis code, and list the claims an adversarial review
overturned.
